\documentclass[fleqn]{article}
\usepackage{amsmath}
\usepackage{amsfonts}
\usepackage{enumitem}
\title{Solución Problemas Sucesiones y Series Funcionales}
\author{Juan Rodríguez}

\setlength{\parindent}{0pt}

\begin{document}
\maketitle
\section{Ejercicio 1}
	Dada la familia de funciones

	\[
	f_n(x)=x^n
	\]

	definidas en el intervalo $[0,+\infty)$, estudiar el campo de convergencia de la sucesión $\{f_n\}$.
	¿Es la sucesión uniformemente convergente en $[0,1]$?

	\subsection{Solución}

	Estudiamos el límite puntual:

	\[
	\lim_{n\to\infty}x^n
	\]

	Si $0\leq x<1$:

	\[
	\lim_{n\to\infty}x^n=0
	\]

	Si $x=1$:

	\[
	\lim_{n\to\infty}1^n=1
	\]

	Si $x>1$:

	\[
	\lim_{n\to\infty}x^n=+\infty
	\]

	Por tanto, el campo de convergencia es:

	\[
	\boxed{C=[0,1]}
	\]

	y la función límite es:

	\[
	f(x)=
	\begin{cases}
	0, & 0\leq x<1 \\[1mm]
	1, & x=1
	\end{cases}
	\]

	Estudiamos ahora la convergencia uniforme en $[0,1]$.

	Cada función $f_n(x)=x^n$ es continua en $[0,1]$, pero la función límite $f(x)$ no es continua en $x=1$.

	Por tanto, la convergencia no puede ser uniforme:

	\[
	\boxed{
	\{f_n\}\text{ no converge uniformemente en }[0,1]
	}
	\]

\section{Ejercicio 2}
	Dada la sucesión de funciones cuyo término general es

	\[
	f_n(x)=x+n
	\]

	estudiar su convergencia.

	\subsection{Solución}

	Estudiamos el límite puntual:

	\[
	\lim_{n\to\infty}f_n(x)
	=
	\lim_{n\to\infty}(x+n)
	\]

	Para cualquier $x\in\mathbb{R}$:

	\[
	\lim_{n\to\infty}(x+n)=+\infty
	\]

	Por tanto, la sucesión no converge puntualmente a ninguna función real.

	\[
	\boxed{
	C=\emptyset
	}
	\]

	En consecuencia, tampoco puede existir convergencia uniforme en ningún intervalo no vacío.

	\[
	\boxed{
	\{f_n\}\text{ no es uniformemente convergente}
	}
	\]

\section{Ejercicio 3}
	Dada la sucesión de funciones cuyo término general es

	\[
	f_n(x)=e^{-nx}
	\]

	determinar su límite puntual y si la sucesión converge uniformemente en el intervalo
	$[a,+\infty)$ con $a>0$.

	\subsection{Solución}

	Estudiamos el límite puntual:

	\[
	\lim_{n\to\infty}e^{-nx}
	\]

	Si $x>0$:

	\[
	\lim_{n\to\infty}e^{-nx}=0
	\]

	Si $x=0$:

	\[
	\lim_{n\to\infty}e^0=1
	\]

	Si $x<0$:

	\[
	\lim_{n\to\infty}e^{-nx}=+\infty
	\]

	Por tanto, el campo de convergencia es:

	\[
	\boxed{
	C=[0,+\infty)
	}
	\]

	y la función límite es:

	\[
	f(x)=
	\begin{cases}
	1, & x=0 \\[1mm]
	0, & x>0
	\end{cases}
	\]

	Estudiamos ahora la convergencia uniforme en $[a,+\infty)$, con $a>0$.

	Como en dicho intervalo $f(x)=0$:

	\[
	|f_n(x)-f(x)|
	=
	e^{-nx}
	\]

	El máximo se alcanza en $x=a$:

	\[
	\sup_{x\in[a,+\infty)}
	|f_n(x)-f(x)|
	=
	e^{-na}
	\]

	\[
	\lim_{n\to\infty}e^{-na}=0
	\]

	Por tanto:

	\[
	\boxed{
	\{f_n\}\text{ converge uniformemente en }[a,+\infty)
	}
	\]

\section{Ejercicio 4}
	Dada la sucesión de funciones cuyo término general es

	\[
	f_n(x)=nxe^{-nx}
	\]

	determinar su límite puntual y si la sucesión converge uniformemente en algún intervalo.

	\subsection{Solución}

	Estudiamos el límite puntual:

	\[
	\lim_{n\to\infty}nxe^{-nx}
	\]

	Si $x>0$:

	\[
	\lim_{n\to\infty}nxe^{-nx}=0
	\]

	Si $x=0$:

	\[
	f_n(0)=0
	\]

	Si $x<0$:

	\[
	nxe^{-nx}\to-\infty
	\]

	Por tanto, el campo de convergencia es:

	\[
	\boxed{
	C=[0,+\infty)
	}
	\]

	y la función límite es:

	\[
	f(x)=0
	\]

	Estudiamos ahora la convergencia uniforme en $[0,+\infty)$:

	\[
	|f_n(x)-f(x)|
	=
	nxe^{-nx}
	\]

	Buscamos su máximo:

	\[
	\frac{d}{dx}
	\left(
	nxe^{-nx}
	\right)
	=
	ne^{-nx}(1-nx)
	\]

	\[
	1-nx=0
	\qquad\Longrightarrow\qquad
	x=\frac{1}{n}
	\]

	Por tanto:

	\[
	\sup_{x\in[0,+\infty)}
	|f_n(x)|
	=
	f_n\left(\frac{1}{n}\right)
	=
	e^{-1}
	\]

	Como:

	\[
	\lim_{n\to\infty}e^{-1}
	=
	e^{-1}\neq0
	\]

	la convergencia no es uniforme en $[0,+\infty)$.

	\[
	\boxed{
	\{f_n\}\text{ no converge uniformemente en }[0,+\infty)
	}
	\]

	Sin embargo, para cualquier $a>0$:

	\[
	\sup_{x\in[a,+\infty)}
	nxe^{-nx}
	=
	nae^{-na}
	\]

	para $n$ suficientemente grande, y:

	\[
	\lim_{n\to\infty}nae^{-na}=0
	\]

	Por tanto:

	\[
	\boxed{
	\{f_n\}\text{ converge uniformemente en }[a,+\infty),
	\quad a>0
	}
	\]

\section{Ejercicio 5}
	Dada la sucesión de funciones cuyo término general es

	\[
	f_n(x)=\frac{2nx}{1+n^4x^2},
	\qquad x\in[0,1]
	\]

	determinar su límite puntual y si la sucesión converge uniformemente en dicho intervalo.

	\subsection{Solución}

	Estudiamos el límite puntual:

	\[
	\lim_{n\to\infty}
	\frac{2nx}{1+n^4x^2}
	\]

	Si $x=0$:

	\[
	f_n(0)=0
	\]

	Si $x>0$:

	\[
	\frac{2nx}{1+n^4x^2}
	=
	\frac{2x/n^3}{1/n^4+x^2}
	\longrightarrow 0
	\]

	Por tanto:

	\[
	\boxed{
	f(x)=0
	\qquad \forall x\in[0,1]
	}
	\]

	Estudiamos ahora la convergencia uniforme:

	\[
	|f_n(x)-f(x)|
	=
	\frac{2nx}{1+n^4x^2}
	\]

	Buscamos su máximo derivando:

	\[
	f_n'(x)
	=
	\frac{
	2n(1+n^4x^2)-2nx(2n^4x)
	}{
	(1+n^4x^2)^2
	}
	\]

	\[
	=
	\frac{
	2n(1-n^4x^2)
	}{
	(1+n^4x^2)^2
	}
	\]

	\[
	1-n^4x^2=0
	\qquad\Longrightarrow\qquad
	x=\frac{1}{n^2}
	\]

	Entonces:

	\[
	\sup_{x\in[0,1]}
	|f_n(x)|
	=
	f_n\left(\frac{1}{n^2}\right)
	=
	\frac{2n\frac{1}{n^2}}{1+n^4\frac{1}{n^4}}
	=
	\frac{1}{n}
	\]

	Como:

	\[
	\lim_{n\to\infty}\frac{1}{n}=0
	\]

	se concluye que:

	\[
	\boxed{
	\{f_n\}\text{ converge uniformemente en }[0,1]
	}
	\]

\section{Ejercicio 6}
	Dada la sucesión de funciones cuyo término general es

	\[
	f_n(x)=\frac{nx}{1+nx^2},
	\qquad x\in[0,+\infty)
	\]

	calcular su límite puntual. A continuación, estudiar si esa función y las $f_n$ están acotadas en ese intervalo. ¿Es la convergencia uniforme?

	\subsection{Solución}

	Estudiamos el límite puntual:

	\[
	\lim_{n\to\infty}
	\frac{nx}{1+nx^2}
	\]

	Si $x=0$:

	\[
	f_n(0)=0
	\]

	Si $x>0$:

	\[
	\frac{nx}{1+nx^2}
	=
	\frac{x}{\frac{1}{n}+x^2}
	\longrightarrow
	\frac{1}{x}
	\]

	Por tanto:

	\[
	f(x)=
	\begin{cases}
	0, & x=0 \\[1mm]
	\displaystyle\frac{1}{x}, & x>0
	\end{cases}
	\]

	La función límite no está acotada en $[0,+\infty)$, ya que:

	\[
	\lim_{x\to0^+}\frac{1}{x}=+\infty
	\]

	En cambio, para cada $n$ fijo, estudiamos el máximo de $f_n$:

	\[
	f_n'(x)
	=
	\frac{
	n(1+nx^2)-nx(2nx)
	}{
	(1+nx^2)^2
	}
	=
	\frac{
	n(1-nx^2)
	}{
	(1+nx^2)^2
	}
	\]

	\[
	1-nx^2=0
	\qquad\Longrightarrow\qquad
	x=\frac{1}{\sqrt{n}}
	\]

	\[
	f_n\left(\frac{1}{\sqrt{n}}\right)
	=
	\frac{\sqrt{n}}{2}
	\]

	Por tanto, cada función $f_n$ está acotada en $[0,+\infty)$.

	Como todas las $f_n$ están acotadas pero la función límite $f$ no está acotada, la convergencia no puede ser uniforme.

	\[
	\boxed{
	\{f_n\}\text{ no converge uniformemente en }[0,+\infty)
	}
	\]

\section{Ejercicio 7}
	Estudiar la convergencia puntual y uniforme de la sucesión de funciones cuyo término general es

	\[
	f_n(x)=\arctan(nx)
	\]

	en todo $\mathbb{R}$.

	\subsection{Solución}

	Estudiamos el límite puntual:

	\[
	\lim_{n\to\infty}\arctan(nx)
	\]

	Si $x>0$:

	\[
	\lim_{n\to\infty}\arctan(nx)
	=
	\frac{\pi}{2}
	\]

	Si $x=0$:

	\[
	f_n(0)=0
	\]

	Si $x<0$:

	\[
	\lim_{n\to\infty}\arctan(nx)
	=
	-\frac{\pi}{2}
	\]

	Por tanto, la función límite es:

	\[
	f(x)=
	\begin{cases}
	-\frac{\pi}{2}, & x<0 \\[1mm]
	0, & x=0 \\[1mm]
	\frac{\pi}{2}, & x>0
	\end{cases}
	\]

	La función límite no es continua en $x=0$, mientras que todas las funciones
	$f_n(x)=\arctan(nx)$ son continuas en $\mathbb{R}$.

	Por tanto, la convergencia no puede ser uniforme en $\mathbb{R}$:

	\[
	\boxed{
	\{f_n\}\text{ no converge uniformemente en }\mathbb{R}
	}
	\]

	En cambio, la sucesión sí converge puntualmente en todo $\mathbb{R}$:

	\[
	\boxed{
	C=\mathbb{R}
	}
	\]

\section{Ejercicio 8}
	Dada la sucesión de funciones cuyo término general es

	\[
	f_n(x)=\frac{x+n}{1+xn}
	\]

	calcular su límite puntual en $(0,+\infty)$ y determinar si es uniformemente convergente en $(1,+\infty)$.

	\subsection{Solución}

	Estudiamos el límite puntual:

	\[
	\lim_{n\to\infty}
	\frac{x+n}{1+xn}
	\]

	Dividimos numerador y denominador entre $n$:

	\[
	\frac{x+n}{1+xn}
	=
	\frac{\frac{x}{n}+1}{\frac{1}{n}+x}
	\]

	Por tanto:

	\[
	\lim_{n\to\infty}
	\frac{x+n}{1+xn}
	=
	\frac{1}{x}
	\]

	Así, la función límite es:

	\[
	\boxed{
	f(x)=\frac{1}{x},
	\qquad x\in(0,+\infty)
	}
	\]

	Estudiamos ahora la convergencia uniforme en $(1,+\infty)$:

	\[
	\left|
	f_n(x)-f(x)
	\right|
	=
	\left|
	\frac{x+n}{1+xn}
	-
	\frac{1}{x}
	\right|
	\]

	\[
	=
	\left|
	\frac{x(x+n)-(1+xn)}
	{x(1+xn)}
	\right|
	=
	\frac{x^2-1}{x(1+xn)}
	\]

	Para $x>1$:

	\[
	0<
	\frac{x^2-1}{x(1+xn)}
	<
	\frac{x^2}{x^2n}
	=
	\frac{1}{n}
	\]

	Por tanto:

	\[
	\sup_{x\in(1,+\infty)}
	|f_n(x)-f(x)|
	\leq
	\frac{1}{n}
	\]

	Como:

	\[
	\lim_{n\to\infty}\frac{1}{n}=0
	\]

	se concluye que:

	\[
	\boxed{
	\{f_n\}\text{ converge uniformemente en }(1,+\infty)
	}
	\]

\section{Ejercicio 9}
	Demostrar que la sucesión $\{f_n\}$, donde

	\[
	f_n(x)=\frac{2n^2x}{(n^2x^2+1)^2}
	\]

	no es uniformemente convergente en el intervalo $[0,1]$ utilizando integrales.

	\subsection{Solución}

	Primero calculamos el límite puntual.

	Si $x=0$:

	\[
	f_n(0)=0
	\]

	Si $x>0$:

	\[
	\frac{2n^2x}{(n^2x^2+1)^2}
	=
	\frac{\frac{2x}{n^2}}
	{\left(x^2+\frac{1}{n^2}\right)^2}
	\longrightarrow 0
	\]

	Por tanto:

	\[
	\boxed{
	f(x)=0
	}
	\]

	Si la convergencia fuese uniforme, al ser las $f_n$ continuas en $[0,1]$ se cumpliría:

	\[
	\lim_{n\to\infty}
	\int_0^1 f_n(x)\,dx
	=
	\int_0^1
	\lim_{n\to\infty}f_n(x)\,dx
	=
	0
	\]

	Calculamos la integral:

	\[
	\int_0^1
	\frac{2n^2x}{(n^2x^2+1)^2}\,dx
	\]

	Hacemos:

	\[
	u=n^2x^2+1,
	\qquad
	du=2n^2x\,dx
	\]

	\[
	\int_0^1
	\frac{2n^2x}{(n^2x^2+1)^2}\,dx
	=
	\int_1^{n^2+1}\frac{1}{u^2}\,du
	\]

	\[
	=
	\left[
	-\frac{1}{u}
	\right]_1^{n^2+1}
	=
	1-\frac{1}{n^2+1}
	\]

	Por tanto:

	\[
	\lim_{n\to\infty}
	\int_0^1f_n(x)\,dx
	=
	1
	\neq
	0
	=
	\int_0^1f(x)\,dx
	\]

	Luego la convergencia no puede ser uniforme:

	\[
	\boxed{
	\{f_n\}\text{ no converge uniformemente en }[0,1]
	}
	\]

\section{Ejercicio 10}
	Estudiar la convergencia puntual de la serie de funciones

	\[
	\sum_{n=1}^{\infty}
	\left(
	x+\frac{1}{n}
	\right)^n,
	\qquad x\geq 0
	\]

	\subsection{Solución}

	Para que la serie pueda converger es necesario que:

	\[
	\lim_{n\to\infty}
	\left(
	x+\frac{1}{n}
	\right)^n
	=
	0
	\]

	Estudiamos según el valor de $x$.

	Si $x=0$:

	\[
	\left(
	\frac{1}{n}
	\right)^n
	\longrightarrow 0
	\]

	y la serie

	\[
	\sum_{n=1}^{\infty}
	\frac{1}{n^n}
	\]

	es convergente.

	Si $0<x<1$, aplicamos el criterio de la raíz:

	\[
	\sqrt[n]{
	\left(
	x+\frac{1}{n}
	\right)^n
	}
	=
	x+\frac{1}{n}
	\]

	\[
	\lim_{n\to\infty}
	\left(
	x+\frac{1}{n}
	\right)
	=
	x<1
	\]

	Por tanto, la serie converge.

	Si $x=1$:

	\[
	\left(
	1+\frac{1}{n}
	\right)^n
	\longrightarrow e\neq0
	\]

	por lo que la serie diverge.

	Si $x>1$:

	\[
	\left(
	x+\frac{1}{n}
	\right)^n
	\not\longrightarrow0
	\]

	y la serie también diverge.

	Por tanto, el campo de convergencia puntual es:

	\[
	\boxed{
	C=[0,1)
	}
	\]

\section{Ejercicio 11}
	Estudiar la convergencia puntual de la serie de funciones

	\[
	\sum_{n=1}^{\infty}x^{n-1}
	\]

	en $(0,1)$ y calcular $S(x)$.

	\subsection{Solución}

	La serie es geométrica:

	\[
	\sum_{n=1}^{\infty}x^{n-1}
	=
	\sum_{k=0}^{\infty}x^k
	\]

	Una serie geométrica converge si:

	\[
	|x|<1
	\]

	Como $x\in(0,1)$, la serie converge para todo punto del intervalo.

	Por tanto:

	\[
	\boxed{
	C=(0,1)
	}
	\]

	La suma de una serie geométrica es:

	\[
	\sum_{k=0}^{\infty}x^k
	=
	\frac{1}{1-x}
	\]

	Por tanto:

	\[
	\boxed{
	S(x)=\frac{1}{1-x},
	\qquad x\in(0,1)
	}
	\]

\section{Ejercicio 12}
	Determinar si la serie de funciones

	\[
	\sum_{n=1}^{\infty}(1-x)x^{n-1}
	\]

	es uniformemente convergente en $[0,1]$.

	\subsection{Solución}

	Para $0\leq x<1$, la serie es geométrica:

	\[
	\sum_{n=1}^{\infty}(1-x)x^{n-1}
	=
	(1-x)\sum_{n=1}^{\infty}x^{n-1}
	\]

	\[
	=
	(1-x)\frac{1}{1-x}
	=
	1
	\]

	Si $x=1$:

	\[
	(1-x)x^{n-1}=0
	\]

	para todo $n$, por lo que:

	\[
	S(1)=0
	\]

	Así, la función suma es:

	\[
	S(x)=
	\begin{cases}
	1, & 0\leq x<1 \\[1mm]
	0, & x=1
	\end{cases}
	\]

	Cada término

	\[
	f_n(x)=(1-x)x^{n-1}
	\]

	es continuo en $[0,1]$, pero la función suma $S(x)$ no es continua en $x=1$.

	Por tanto, la serie no puede converger uniformemente en $[0,1]$:

	\[
	\boxed{
	\sum_{n=1}^{\infty}(1-x)x^{n-1}
	\text{ no converge uniformemente en }[0,1]
	}
	\]

\section{Ejercicio 13}
	Determinar si la serie de funciones

	\[
	\sum_{n=1}^{\infty}\frac{\sin(nx)}{n^2}
	\]

	es uniformemente convergente en todo $\mathbb{R}$.

	\subsection{Solución}

	Aplicamos el criterio mayorante de Weierstrass.

	Como:

	\[
	|\sin(nx)|\leq 1
	\]

	se tiene:

	\[
	\left|
	\frac{\sin(nx)}{n^2}
	\right|
	\leq
	\frac{1}{n^2}
	\]

	para todo $x\in\mathbb{R}$.

	La serie numérica:

	\[
	\sum_{n=1}^{\infty}\frac{1}{n^2}
	\]

	es convergente.

	Por el criterio mayorante de Weierstrass:

	\[
	\boxed{
	\sum_{n=1}^{\infty}\frac{\sin(nx)}{n^2}
	\text{ converge uniformemente en }\mathbb{R}
	}
	\]

\section{Ejercicio 14}
	Determinar si la serie de funciones

	\[
	\sum_{n=1}^{\infty}
	\frac{x}{1+n^2x^2}
	\]

	es uniformemente convergente en el intervalo $[c,+\infty)$, donde $c>0$.

	\subsection{Solución}

	Aplicamos el criterio mayorante de Weierstrass.

	Para $x\geq c>0$:

	\[
	\frac{x}{1+n^2x^2}
	\leq
	\frac{x}{n^2x^2}
	=
	\frac{1}{n^2x}
	\]

	Como $x\geq c$:

	\[
	\frac{1}{n^2x}
	\leq
	\frac{1}{n^2c}
	\]

	Por tanto:

	\[
	\left|
	\frac{x}{1+n^2x^2}
	\right|
	\leq
	\frac{1}{n^2c}
	\]

	La serie numérica:

	\[
	\sum_{n=1}^{\infty}
	\frac{1}{n^2c}
	=
	\frac{1}{c}
	\sum_{n=1}^{\infty}
	\frac{1}{n^2}
	\]

	es convergente.

	Por el criterio mayorante de Weierstrass:

	\[
	\boxed{
	\sum_{n=1}^{\infty}
	\frac{x}{1+n^2x^2}
	\text{ converge uniformemente en }[c,+\infty)
	}
	\]

\section{Ejercicio 15}
	Determinar si la serie de funciones

	\[
	\sum_{n=0}^{\infty}
	\frac{(-1)^n}{n!}x^n
	\]

	es uniformemente convergente en el intervalo $[-a,a]$, donde $a>0$.

	\subsection{Solución}

	Aplicamos el criterio mayorante de Weierstrass.

	Para $x\in[-a,a]$:

	\[
	|x|\leq a
	\]

	Por tanto:

	\[
	\left|
	\frac{(-1)^n}{n!}x^n
	\right|
	=
	\frac{|x|^n}{n!}
	\leq
	\frac{a^n}{n!}
	\]

	La serie numérica:

	\[
	\sum_{n=0}^{\infty}\frac{a^n}{n!}
	\]

	es convergente, ya que:

	\[
	\sum_{n=0}^{\infty}\frac{a^n}{n!}
	=
	e^a
	\]

	Por el criterio mayorante de Weierstrass:

	\[
	\boxed{
	\sum_{n=0}^{\infty}
	\frac{(-1)^n}{n!}x^n
	\text{ converge uniformemente en }[-a,a]
	}
	\]

\section{Ejercicio 16}
	Demostrar que la serie de funciones

	\[
	\sum_{n=0}^{\infty}
	e^{-n^2\alpha}\sin(nx)
	\]

	es uniformemente convergente en $\mathbb{R}$ para cualquier valor $\alpha>0$.

	\subsection{Solución}

	Aplicamos el criterio mayorante de Weierstrass.

	Como:

	\[
	|\sin(nx)|\leq 1
	\]

	se tiene:

	\[
	\left|
	e^{-n^2\alpha}\sin(nx)
	\right|
	\leq
	e^{-n^2\alpha}
	\]

	Para $n\geq1$:

	\[
	n^2\geq n
	\]

	y como $\alpha>0$:

	\[
	e^{-n^2\alpha}
	\leq
	e^{-n\alpha}
	\]

	La serie:

	\[
	\sum_{n=1}^{\infty}e^{-n\alpha}
	\]

	es geométrica de razón:

	\[
	e^{-\alpha}<1
	\]

	por lo que es convergente.

	Por comparación:

	\[
	\sum_{n=0}^{\infty}e^{-n^2\alpha}
	\]

	es convergente.

	Por el criterio mayorante de Weierstrass:

	\[
	\boxed{
	\sum_{n=0}^{\infty}
	e^{-n^2\alpha}\sin(nx)
	\text{ converge uniformemente en }\mathbb{R}
	}
	\]

\section{Ejercicio 17}
	Determinar si la serie

	\[
	\sum_{n=1}^{\infty}
	\frac{n^2+x^4}{n^4+x^2}
	\]

	converge uniformemente a una cierta función $S(x)$ en $[-a,a]$, donde $a>0$.
	¿Se puede afirmar que $S(x)$ es continua en $[-a,a]$?

	\subsection{Solución}

	Aplicamos el criterio mayorante de Weierstrass.

	Para $x\in[-a,a]$:

	\[
	|x|\leq a
	\]

	y por tanto:

	\[
	x^4\leq a^4
	\]

	Además:

	\[
	n^4+x^2\geq n^4
	\]

	Así:

	\[
	0\leq
	\frac{n^2+x^4}{n^4+x^2}
	\leq
	\frac{n^2+a^4}{n^4}
	\]

	\[
	=
	\frac{1}{n^2}
	+
	\frac{a^4}{n^4}
	\]

	La serie numérica:

	\[
	\sum_{n=1}^{\infty}
	\left(
	\frac{1}{n^2}
	+
	\frac{a^4}{n^4}
	\right)
	\]

	es convergente.

	Por el criterio mayorante de Weierstrass:

	\[
	\boxed{
	\sum_{n=1}^{\infty}
	\frac{n^2+x^4}{n^4+x^2}
	\text{ converge uniformemente en }[-a,a]
	}
	\]

	Cada función:

	\[
	f_n(x)=
	\frac{n^2+x^4}{n^4+x^2}
	\]

	es continua en $[-a,a]$.

	Como la serie converge uniformemente y cada $f_n$ es continua, su función suma $S(x)$ también es continua.

	\[
	\boxed{
	S(x)\text{ es continua en }[-a,a]
	}
	\]

\section{Ejercicio 18}
	Calcular el campo de convergencia de la serie de potencias

	\[
	\sum_{n=0}^{\infty}x^n
	\]

	\subsection{Solución}

	Aplicamos el criterio del cociente:

	\[
	\lim_{n\to\infty}
	\left|
	\frac{x^{n+1}}{x^n}
	\right|
	=
	|x|
	\]

	La serie converge si:

	\[
	|x|<1
	\]

	Por tanto, el radio de convergencia es:

	\[
	\boxed{
	R=1
	}
	\]

	Estudiamos los extremos.

	Si $x=1$:

	\[
	\sum_{n=0}^{\infty}1
	\]

	es divergente.

	Si $x=-1$:

	\[
	\sum_{n=0}^{\infty}(-1)^n
	\]

	es divergente, ya que el término general no tiende a cero.

	Por tanto, el campo de convergencia es:

	\[
	\boxed{
	C=(-1,1)
	}
	\]

\section{Ejercicio 19}
	Calcular el campo de convergencia de la serie de potencias

	\[
	\sum_{n=1}^{\infty}n!x^n
	\]

	\subsection{Solución}

	Aplicamos el criterio del cociente:

	\[
	\lim_{n\to\infty}
	\left|
	\frac{(n+1)!x^{n+1}}{n!x^n}
	\right|
	\]

	\[
	=
	\lim_{n\to\infty}
	(n+1)|x|
	\]

	Si $x\neq0$:

	\[
	\lim_{n\to\infty}
	(n+1)|x|
	=
	+\infty
	\]

	por lo que la serie diverge.

	Si $x=0$:

	\[
	\sum_{n=1}^{\infty}n!0^n=0
	\]

	y la serie converge.

	Por tanto, el radio de convergencia es:

	\[
	\boxed{
	R=0
	}
	\]

	y el campo de convergencia es:

	\[
	\boxed{
	C=\{0\}
	}
	\]

\section{Ejercicio 20}
	Calcular el campo de convergencia de la serie de potencias

	\[
	\sum_{n=1}^{\infty}\frac{x^n}{n!}
	\]

	\subsection{Solución}

	Aplicamos el criterio del cociente:

	\[
	\lim_{n\to\infty}
	\left|
	\frac{x^{n+1}}{(n+1)!}
	\frac{n!}{x^n}
	\right|
	\]

	\[
	=
	\lim_{n\to\infty}
	\frac{|x|}{n+1}
	=
	0
	\]

	para cualquier $x\in\mathbb{R}$.

	Como:

	\[
	0<1
	\]

	la serie converge para todo $x\in\mathbb{R}$.

	Por tanto, el radio de convergencia es:

	\[
	\boxed{
	R=+\infty
	}
	\]

	y el campo de convergencia es:

	\[
	\boxed{
	C=\mathbb{R}
	}
	\]

\section{Ejercicio 21}
	Calcular el campo de convergencia de la serie de potencias

	\[
	\sum_{n=0}^{\infty}\frac{x^n}{2n+1}
	\]

	\subsection{Solución}

	Aplicamos el criterio del cociente:

	\[
	\lim_{n\to\infty}
	\left|
	\frac{x^{n+1}}{2n+3}
	\frac{2n+1}{x^n}
	\right|
	\]

	\[
	=
	|x|
	\lim_{n\to\infty}
	\frac{2n+1}{2n+3}
	=
	|x|
	\]

	La serie converge si:

	\[
	|x|<1
	\]

	Por tanto, el radio de convergencia es:

	\[
	\boxed{
	R=1
	}
	\]

	Estudiamos los extremos.

	Si $x=1$:

	\[
	\sum_{n=0}^{\infty}\frac{1}{2n+1}
	\]

	es divergente.

	Si $x=-1$:

	\[
	\sum_{n=0}^{\infty}\frac{(-1)^n}{2n+1}
	\]

	es convergente por el criterio de Leibniz.

	Por tanto, el campo de convergencia es:

	\[
	\boxed{
	C=[-1,1)
	}
	\]

\section{Ejercicio 22}
	Calcular el campo de convergencia de la serie de potencias

	\[
	\sum_{n=1}^{\infty}2^n x^n
	\]

	\subsection{Solución}

	Aplicamos el criterio del cociente:

	\[
	\lim_{n\to\infty}
	\left|
	\frac{2^{n+1}x^{n+1}}{2^n x^n}
	\right|
	=
	2|x|
	\]

	La serie converge si:

	\[
	2|x|<1
	\]

	\[
	|x|<\frac{1}{2}
	\]

	Por tanto, el radio de convergencia es:

	\[
	\boxed{
	R=\frac{1}{2}
	}
	\]

	Estudiamos los extremos.

	Si $x=\frac{1}{2}$:

	\[
	\sum_{n=1}^{\infty}
	2^n\left(\frac{1}{2}\right)^n
	=
	\sum_{n=1}^{\infty}1
	\]

	es divergente.

	Si $x=-\frac{1}{2}$:

	\[
	\sum_{n=1}^{\infty}
	2^n\left(-\frac{1}{2}\right)^n
	=
	\sum_{n=1}^{\infty}(-1)^n
	\]

	es divergente.

	Por tanto, el campo de convergencia es:

	\[
	\boxed{
	C=\left(-\frac{1}{2},\frac{1}{2}\right)
	}
	\]

\section{Ejercicio 23}
	Calcular el campo de convergencia de la serie de potencias

	\[
	\sum_{n=1}^{\infty}\frac{x^{2n}}{3^n}
	\]

	\subsection{Solución}

	Podemos escribir:

	\[
	\frac{x^{2n}}{3^n}
	=
	\left(
	\frac{x^2}{3}
	\right)^n
	\]

	Por tanto, es una serie geométrica de razón:

	\[
	r=\frac{x^2}{3}
	\]

	La serie converge si:

	\[
	\left|
	\frac{x^2}{3}
	\right|<1
	\]

	Como $x^2\geq0$:

	\[
	\frac{x^2}{3}<1
	\]

	\[
	x^2<3
	\]

	\[
	|x|<\sqrt{3}
	\]

	Por tanto, el radio de convergencia es:

	\[
	\boxed{
	R=\sqrt{3}
	}
	\]

	Estudiamos los extremos.

	Si $x=\sqrt{3}$:

	\[
	\sum_{n=1}^{\infty}
	\frac{(\sqrt{3})^{2n}}{3^n}
	=
	\sum_{n=1}^{\infty}1
	\]

	es divergente.

	Si $x=-\sqrt{3}$:

	\[
	\sum_{n=1}^{\infty}
	\frac{(-\sqrt{3})^{2n}}{3^n}
	=
	\sum_{n=1}^{\infty}1
	\]

	es divergente.

	Por tanto, el campo de convergencia es:

	\[
	\boxed{
	C=(-\sqrt{3},\sqrt{3})
	}
	\]

\section{Ejercicio 24}
	Calcular el campo de convergencia de la serie de potencias

	\[
	\sum_{n=1}^{\infty}
	\frac{(x-2)^n}{n^2}
	\]

	\subsection{Solución}

	Aplicamos el criterio del cociente:

	\[
	\lim_{n\to\infty}
	\left|
	\frac{(x-2)^{n+1}}{(n+1)^2}
	\frac{n^2}{(x-2)^n}
	\right|
	\]

	\[
	=
	|x-2|
	\lim_{n\to\infty}
	\frac{n^2}{(n+1)^2}
	=
	|x-2|
	\]

	La serie converge si:

	\[
	|x-2|<1
	\]

	\[
	1<x<3
	\]

	Por tanto, el radio de convergencia es:

	\[
	\boxed{
	R=1
	}
	\]

	Estudiamos los extremos.

	Si $x=1$:

	\[
	\sum_{n=1}^{\infty}
	\frac{(-1)^n}{n^2}
	\]

	es convergente.

	Si $x=3$:

	\[
	\sum_{n=1}^{\infty}
	\frac{1}{n^2}
	\]

	es convergente.

	Por tanto, el campo de convergencia es:

	\[
	\boxed{
	C=[1,3]
	}
	\]

\section{Ejercicio 25}
	Determinar el desarrollo en serie de la función

	\[
	f(x)=(1+e^x)^2
	\]

	y calcular su radio de convergencia.

	\subsection{Solución}

	Desarrollamos:

	\[
	(1+e^x)^2
	=
	1+2e^x+e^{2x}
	\]

	Sabemos que:

	\[
	e^x
	=
	\sum_{n=0}^{\infty}\frac{x^n}{n!}
	\]

	y:

	\[
	e^{2x}
	=
	\sum_{n=0}^{\infty}\frac{2^n x^n}{n!}
	\]

	Por tanto:

	\[
	f(x)
	=
	1
	+
	2\sum_{n=0}^{\infty}\frac{x^n}{n!}
	+
	\sum_{n=0}^{\infty}\frac{2^n x^n}{n!}
	\]

	Separando el término constante:

	\[
	f(x)
	=
	4
	+
	\sum_{n=1}^{\infty}
	\frac{2+2^n}{n!}x^n
	\]

	Así, el desarrollo en serie de potencias es:

	\[
	\boxed{
	(1+e^x)^2
	=
	4+
	\sum_{n=1}^{\infty}
	\frac{2+2^n}{n!}x^n
	}
	\]

	Como las series de $e^x$ y $e^{2x}$ convergen para todo $x\in\mathbb{R}$:

	\[
	\boxed{
	R=+\infty
	}
	\]

\section{Ejercicio 26}
	Determinar el desarrollo en serie de la función

	\[
	f(x)=\frac{1}{\sqrt{2-x}}
	\]

	y calcular su radio de convergencia.

	\subsection{Solución}

	Escribimos:

	\[
	\frac{1}{\sqrt{2-x}}
	=
	\frac{1}{\sqrt{2}}
	\frac{1}{\sqrt{1-\frac{x}{2}}}
	\]

	Usamos el desarrollo:

	\[
	(1-u)^{-1/2}
	=
	\sum_{n=0}^{\infty}
	\frac{(2n)!}{4^n(n!)^2}u^n
	\]

	con:

	\[
	u=\frac{x}{2}
	\]

	Por tanto:

	\[
	\frac{1}{\sqrt{2-x}}
	=
	\frac{1}{\sqrt{2}}
	\sum_{n=0}^{\infty}
	\frac{(2n)!}{4^n(n!)^2}
	\left(
	\frac{x}{2}
	\right)^n
	\]

	\[
	=
	\frac{1}{\sqrt{2}}
	\sum_{n=0}^{\infty}
	\frac{(2n)!}{8^n(n!)^2}x^n
	\]

	Por tanto:

	\[
	\boxed{
	\frac{1}{\sqrt{2-x}}
	=
	\frac{1}{\sqrt{2}}
	\sum_{n=0}^{\infty}
	\frac{(2n)!}{8^n(n!)^2}x^n
	}
	\]

	El desarrollo de $(1-u)^{-1/2}$ converge para:

	\[
	|u|<1
	\]

	Así:

	\[
	\left|
	\frac{x}{2}
	\right|<1
	\]

	\[
	|x|<2
	\]

	Por tanto, el radio de convergencia es:

	\[
	\boxed{
	R=2
	}
	\]

\section{Ejercicio 27}
	Hallar el desarrollo en serie de la función

	\[
	f(x)=\frac{3x+2}{x^2-5x+6}
	\]

	y calcular su radio de convergencia.

	\subsection{Solución}

	Factorizamos el denominador:

	\[
	x^2-5x+6=(x-2)(x-3)
	\]

	Descomponemos en fracciones simples:

	\[
	\frac{3x+2}{(x-2)(x-3)}
	=
	\frac{A}{x-2}
	+
	\frac{B}{x-3}
	\]

	\[
	3x+2
	=
	A(x-3)+B(x-2)
	\]

	\[
	A+B=3
	\]

	\[
	-3A-2B=2
	\]

	De donde:

	\[
	A=-8,
	\qquad
	B=11
	\]

	Por tanto:

	\[
	f(x)
	=
	-\frac{8}{x-2}
	+
	\frac{11}{x-3}
	\]

	Reescribimos:

	\[
	-\frac{8}{x-2}
	=
	\frac{4}{1-\frac{x}{2}}
	\]

	y:

	\[
	\frac{11}{x-3}
	=
	-\frac{11}{3}
	\frac{1}{1-\frac{x}{3}}
	\]

	Usando la serie geométrica:

	\[
	\frac{1}{1-u}
	=
	\sum_{n=0}^{\infty}u^n
	\]

	obtenemos:

	\[
	f(x)
	=
	4\sum_{n=0}^{\infty}
	\left(\frac{x}{2}\right)^n
	-
	\frac{11}{3}
	\sum_{n=0}^{\infty}
	\left(\frac{x}{3}\right)^n
	\]

	\[
	=
	\sum_{n=0}^{\infty}
	\left(
	\frac{4}{2^n}
	-
	\frac{11}{3^{n+1}}
	\right)x^n
	\]

	Por tanto:

	\[
	\boxed{
	\frac{3x+2}{x^2-5x+6}
	=
	\sum_{n=0}^{\infty}
	\left(
	\frac{4}{2^n}
	-
	\frac{11}{3^{n+1}}
	\right)x^n
	}
	\]

	La primera serie requiere:

	\[
	\left|\frac{x}{2}\right|<1
	\qquad\Longrightarrow\qquad
	|x|<2
	\]

	y la segunda:

	\[
	\left|\frac{x}{3}\right|<1
	\qquad\Longrightarrow\qquad
	|x|<3
	\]

	Por tanto, el desarrollo completo converge para:

	\[
	|x|<2
	\]

	y el radio de convergencia es:

	\[
	\boxed{
	R=2
	}
	\]

\section{Ejercicio 28}
	Determinar el radio de convergencia de la serie

	\[
	\sum_{n=1}^{\infty}n^2x^n
	\]

	y calcular su suma derivando la expresión de la serie geométrica.

	\subsection{Solución}

	Partimos de la serie geométrica:

	\[
	\sum_{n=0}^{\infty}x^n
	=
	\frac{1}{1-x},
	\qquad |x|<1
	\]

	Derivamos:

	\[
	\sum_{n=1}^{\infty}nx^{n-1}
	=
	\frac{1}{(1-x)^2}
	\]

	Multiplicamos por $x$:

	\[
	\sum_{n=1}^{\infty}nx^n
	=
	\frac{x}{(1-x)^2}
	\]

	Derivamos de nuevo:

	\[
	\sum_{n=1}^{\infty}n^2x^{n-1}
	=
	\frac{d}{dx}
	\left(
	\frac{x}{(1-x)^2}
	\right)
	\]

	\[
	=
	\frac{1+x}{(1-x)^3}
	\]

	Multiplicando por $x$:

	\[
	\boxed{
	\sum_{n=1}^{\infty}n^2x^n
	=
	\frac{x(1+x)}{(1-x)^3}
	}
	\]

	Como la serie geométrica y sus derivadas tienen el mismo radio de convergencia:

	\[
	\boxed{
	R=1
	}
	\]

	La expresión de la suma es válida para:

	\[
	\boxed{
	|x|<1
	}
	\]

\section{Ejercicio 29}
	Dada la serie

	\[
	\sum_{n=1}^{\infty}
	\frac{x^n}{3^{n+1}n!}
	\]

	calcular su intervalo de convergencia y sumar la serie utilizando el desarrollo como serie de potencias de la función $e^x$.

	\subsection{Solución}

	Escribimos:

	\[
	\frac{x^n}{3^{n+1}n!}
	=
	\frac{1}{3}
	\frac{\left(\frac{x}{3}\right)^n}{n!}
	\]

	Por tanto:

	\[
	\sum_{n=1}^{\infty}
	\frac{x^n}{3^{n+1}n!}
	=
	\frac{1}{3}
	\sum_{n=1}^{\infty}
	\frac{\left(\frac{x}{3}\right)^n}{n!}
	\]

	Sabemos que:

	\[
	e^u
	=
	\sum_{n=0}^{\infty}
	\frac{u^n}{n!}
	\]

	Así:

	\[
	\sum_{n=1}^{\infty}
	\frac{\left(\frac{x}{3}\right)^n}{n!}
	=
	e^{x/3}-1
	\]

	Por tanto:

	\[
	\boxed{
	\sum_{n=1}^{\infty}
	\frac{x^n}{3^{n+1}n!}
	=
	\frac{1}{3}
	\left(
	e^{x/3}-1
	\right)
	}
	\]

	Como la serie de $e^x$ converge para todo $x\in\mathbb{R}$, también lo hace la serie dada.

	Por tanto:

	\[
	\boxed{
	R=+\infty
	}
	\]

	y el intervalo de convergencia es:

	\[
	\boxed{
	C=\mathbb{R}
	}
	\]

\section{Ejercicio 30}
	Dada la serie

	\[
	\sum_{n=1}^{\infty}x^n(1-x^n)
	\]

	determinar el campo de convergencia y calcular su suma utilizando las expresiones de otras series de potencias.

	\subsection{Solución}

	Desarrollamos:

	\[
	x^n(1-x^n)
	=
	x^n-x^{2n}
	\]

	Por tanto:

	\[
	\sum_{n=1}^{\infty}x^n(1-x^n)
	=
	\sum_{n=1}^{\infty}x^n
	-
	\sum_{n=1}^{\infty}x^{2n}
	\]

	Ambas son series geométricas.

	La primera converge si:

	\[
	|x|<1
	\]

	y la segunda si:

	\[
	|x^2|<1
	\qquad\Longrightarrow\qquad
	|x|<1
	\]

	Por tanto, para $|x|<1$:

	\[
	\sum_{n=1}^{\infty}x^n
	=
	\frac{x}{1-x}
	\]

	y:

	\[
	\sum_{n=1}^{\infty}x^{2n}
	=
	\frac{x^2}{1-x^2}
	\]

	Así:

	\[
	S(x)
	=
	\frac{x}{1-x}
	-
	\frac{x^2}{1-x^2}
	\]

	\[
	=
	\frac{x(1+x)-x^2}{1-x^2}
	\]

	\[
	=
	\frac{x}{1-x^2}
	\]

	Por tanto:

	\[
	\boxed{
	S(x)=\frac{x}{1-x^2}
	}
	\]

	para:

	\[
	|x|<1
	\]

	Estudiamos los extremos.

	Si $x=1$:

	\[
	x^n(1-x^n)=0
	\]

	para todo $n$, por lo que la serie converge y:

	\[
	S(1)=0
	\]

	Si $x=-1$:

	\[
	(-1)^n\left(1-(-1)^n\right)
	\]

	no tiende a cero, por lo que la serie diverge.

	Por tanto, el campo de convergencia es:

	\[
	\boxed{
	C=(-1,1]
	}
	\]


\end{document}